\documentclass[exercice]{thedocument}%
\usepackage{texmaths-tikz}%
\startdatakeys
\source{}
\cycle{}
\competencesDuSocle{}
\connaissancesRequises{}
\competencesTravaillees{}
\enddatakeys
\begin{document}%
\begin{flushleft}
    Calculer une valeur approchée du périmètre de la figure représentée ci-dessous.
\end{flushleft}
\begin{center}
    \begin{tikzpicture}[scale=1.5]
        \coordinate(A1)at(1,1);
        \coordinate(A2)at(-1,1);
        \coordinate(A3)at(-3,1);
        \coordinate(A4)at(-1,-1);
        \coordinate(A5)at(1,-1);
        \coordinate(A6)at(3,1);
        \draw[fill=Rainbow1-vert](A1)--(A2)--(A3)arc(180:270:2)--(A5)arc(-90:0:2)--cycle;
        \draw[dashed](A5)--(A1);
        \draw[dashed](A4)--(A2);
        \path(A1)--(A2)node[pos=0.5,sloped]{\color{red}||};
        \path(A2)--(A3)node[pos=0.5,sloped]{\color{red}||};
        \path(A2)--(A4)node[pos=0.5,sloped]{\color{red}||};
        \path(A4)--(A5)node[pos=0.5,sloped]{\color{red}||};
        \path(A1)--(A5)node[pos=0.5,sloped]{\color{red}||};
        \path(A1)--(A6)node[pos=0.5,sloped]{\color{red}||};
        \path(A1)--+(0,0.2)coordinate(A1');
        \path(A2)--+(0,0.2)coordinate(A2');
        \draw[latex-latex](A1')--(A2')node[pos=0.5,above]{1,2 dm};
    \end{tikzpicture}
\end{center}
\end{document}
